% ==============================================================================
% T0-Theorie: Dokument 162
% Titel: T0-Optimierungsstrategien: QM-Optimierung und Coherent Ising Machine im Vergleich
% Autor: Johann Pascher
% Datum: Februar 2026
% ==============================================================================

\section*{Abstract}

Die T0-Theorie enthält zwei komplementäre Optimierungsrahmenwerke: Das Dokument~034
(\emph{QM-Optimierung}) beschreibt, wie Quantencomputer durch geometrische Transformation
im zylindrischen Phasenraum und durch Rauschunterdrückung via Zeitfeld-Modulation
optimiert werden können. Das vorliegende Dokument~161 (\emph{T0-Ising Machine}) zeigt,
dass kombinatorische NP-harte Optimierung durch geometrische Relaxation des Massenfeldes
$\mfield$ gelöst werden kann. Dieses Dokument~162 stellt beide Ansätze systematisch
gegenüber und zeigt: Sie sind nicht konkurrierend, sondern \emph{komplementäre Ebenen
desselben geometrischen Prinzips} -- Optimierung als physikalische Relaxation in ein
geometrisches Minimum.

\clearpage

\section{Einführung: Zwei Optimierungsebenen in T0}
\label{sec:intro}

\subsection{Das Ausgangsproblem}

Quantencomputing kämpft auf zwei grundlegend verschiedenen Ebenen mit Rauschen und
Suboptimalität:

\begin{enumerate}
    \item \textbf{Gate-Ebene}: Einzelne Quantengatter sind durch Dekohärenz, Phasenrauschen
    und fraktale Dämpfung gestört. Die Frage: Wie optimiert man \emph{einzelne Operationen}?

    \item \textbf{Problem-Ebene}: Kombinatorische Optimierungsprobleme (NP-hart) sind
    selbst mit perfekten Gattern schwer lösbar. Die Frage: Wie findet man
    \emph{globale Minima} in riesigen Lösungsräumen?
\end{enumerate}

Die T0-Theorie adressiert \textbf{beide Ebenen} -- aber mit unterschiedlichen
Werkzeugen und auf unterschiedlichen Abstraktionsebenen. Dokument~034 \cite{pascher:qmopt}
löst Problem~1; Dokument~161 \cite{pascher:ising} löst Problem~2.

\subsection{Die gemeinsame Wurzel}

Beide Ansätze teilen denselben geometrischen Ursprung: den Parameter
$\xipar = 4/30000 \approx 1{,}333 \times 10^{-4}$
und die Zeitfeld-Bindungsbedingung $T \cdot m = 1$.
Die Unterschiede sind nicht konzeptuell -- sie sind Unterschiede in der
\emph{Skalierungsebene} der Anwendung.

\clearpage

\section{Dokument 034: QM-Optimierung durch geometrischen Formalismus}
\label{sec:doc034}

\subsection{Der zylindrische Phasenraum}

In Dokument~034 \cite{pascher:qmopt} wird der Hilbertraum durch einen
\textbf{zylindrischen Phasenraum} ersetzt. Ein Qubit ist kein 2D-komplexer Vektor,
sondern ein Punkt $(z, r, \theta)$:

\begin{itemize}
    \item $z \in [-1, 1]$: klassische Basisprojektion ($z=+1 \equiv |0\rangle$,
    $z=-1 \equiv |1\rangle$)
    \item $r \geq 0$: Kohärenzradius (Superpositionsgröße), mit $z^2 + r^2 = 1$
    \item $\theta \in [0, 2\pi)$: relative Phase der Überlagerung
\end{itemize}

Gatter-Operationen sind geometrische Transformationen dieses Raumes, keine Matrizenmultiplikationen.

\subsection{Fraktale Dämpfung als intrinsisches Rauschen}

Das X-Gatter (Bit-Flip) ist im T0-Formalismus keine ideale $\pi$-Rotation, sondern
eine gedämpfte Rotation um den Winkel $\alpha = \pi \cdot \Kfrak$:

\begin{align}
    z' &= z \cos(\alpha) - r \sin(\alpha) \\
    r' &= z \sin(\alpha) + r \cos(\alpha)
\end{align}

mit der fraktalen Dämpfungskonstante $\Kfrak = 1 - 100\xipar$.
Ein perfekter Flip ($\alpha = \pi$) ist physikalisch unmöglich -- die Raumzeit
dämpft jede Rotation inhärent. \textbf{Dies ist eine testbare Vorhersage.}

\subsection{Die drei Optimierungsstrategien aus Dok. 034}

\begin{enumerate}
    \item \textbf{T0-Topologie-Compiler}: Qubits werden nicht angeordnet, um SWAP-Operationen
    zu minimieren, sondern um die kumulative \emph{fraktale Weglänge} aller
    Verschränkungsoperationen zu minimieren. Kritisch interagierende Qubits müssen
    physisch näher zusammengebracht werden.

    \item \textbf{Harmonische Resonanzfrequenzen}: Qubit-Frequenzen werden auf
    ''magische'' Frequenzen kalibriert, die sich aus der harmonischen Struktur
    des T0-Zeitfeldes ergeben:
    \begin{equation}
        f_n = \left(\frac{\Ezero}{h}\right) \cdot \xipar^2 \cdot (\phiT^2)^{-n}
        \label{eq:harmonic_freq}
    \end{equation}
    Für supraleitende Qubits ergeben sich Sweet Spots bei ca.\ \textbf{6,24 GHz}
    ($n=14$) und \textbf{2,38 GHz} ($n=15$). Diese Kalibrierung reduziert
    Phasenrauschen intrinsisch.

    \item \textbf{Aktive Zeitfeld-Modulation}: Eine hochfrequente
    \emph{Zeitfeld-Pumpe} bestrahlt untätige Qubits, um das fundamentale
    $\xipar$-Rauschen auszumitteln und die passive $T_2$-Kohärenzgrenze zu
    überwinden -- aktive Dekohärenz-Unterdrückung statt passivem Warten.
\end{enumerate}

\clearpage

\section{Dokument 161: T0-Ising Machine -- Kombinatorische Optimierung}
\label{sec:doc161}

\subsection{Das Ising-Modell als Optimierungssprache}

In Dokument~161 \cite{pascher:ising} wird gezeigt, dass NP-harte kombinatorische
Probleme als Minimierung des Ising-Hamiltonians formuliert werden können:

\begin{equation}
    H_{\text{Ising}} = -\frac{1}{2} \sum_{i \neq j} J_{ij} \, \sigma_i \sigma_j
    - \sum_i h_i \, \sigma_i
    \label{eq:ising}
\end{equation}

Eine Coherent Ising Machine (CIM) findet den Grundzustand dieses Hamiltonians durch
physikalische Relaxation von Laserpulsen (Spins = Phasenzustände 0° oder 180°).

\subsection{Die T0-Ising Machine: Emergente Kopplungen}

Die T0-Theorie erweitert die CIM fundamental: Die Kopplungen $J_{ij}$ sind nicht
extern programmierbar -- sie emergieren geometrisch aus dem Zeitfeld:

\begin{equation}
    J_{ij}^{(\text{T0})} = \int \Tfield_i \cdot \Tfield_j \cdot G(x_i, x_j) \, d^3x
    \label{eq:t0_coupling}
\end{equation}

Das Energiefunktional ist:
\begin{equation}
    H_{\text{Ising}} \leftrightarrow E_{\text{T0}} = \int |\nabla \Tfield|^2 \, d^4x
    \label{eq:energy_equiv}
\end{equation}

\subsection{Rauschen in der CIM: Das Annealing-Problem}

In der CIM spielt Rauschen eine paradoxe Rolle: Es ist gleichzeitig Problem und
Werkzeug. Zu wenig thermisches Rauschen führt zu lokalen Minima-Fallen;
zu viel Rauschen verhindert Konvergenz. Der Annealing-Temperaturparameter
$T_{\text{ann}}$ balanciert diesen Konflikt.

In der T0-Ising Machine übernimmt $\Kfrak^{3/2}$ diese Rolle -- nicht als
externer Parameter, sondern als geometrisch bestimmte fraktale Korrekturgröße.

\clearpage

\section{Systematischer Vergleich: Dok. 034 vs. Dok. 161}
\label{sec:comparison}

\subsection{Übersichtstabelle}

\begin{table}[h]
    \centering
    \caption{Direktvergleich der beiden T0-Optimierungsrahmenwerke}
    \label{tab:main_comparison}
    \resizebox{\textwidth}{!}{\begin{tabular}{@{}llll@{}}
        \toprule
        Aspekt & Dok.\ 034: QM-Optimierung & Dok.\ 161: T0-Ising Machine & Gemeinsame Wurzel \\
        \midrule
        Zielproblem & Gate-Rauschen, Dekohärenz & NP-harte Optimierung & Geometrische Relaxation \\
        Skalierungsebene & Einzelnes Qubit / Gate & Gesamtsystem ($N$ Spins) & $\xipar = 4/30000$ \\
        Spin-Darstellung & $(z, r, \theta)$ im Zylinder & Phase $0°/180°$ (CIM) & $T \cdot m = 1$-Dualität \\
        Rauschrolle & Störung (zu minimieren) & Werkzeug (Annealing) & $\Kfrak$-Korrektur \\
        Kopplung & Gatter-Pulse (extern) & $J_{ij}$ (geometrisch emergiert) & Zeitfeld $\Tfield$ \\
        Optimierungsziel & Maximale Gate-Treue & Globales Minimum von $H$ & Energieminimierung \\
        Determinismus & Explizit (geometr. Trafo) & Versteckt (OPO-Gleichungen) & Beide deterministisch \\
        Testbare Vorhersage & Gedämpfter X-Flip & $\ln(N)$-Residuum nach CIM & $\xipar$-Signaturen \\
        \bottomrule
    \end{tabular}}
\end{table}

\subsection{Parallele 1: Rauschen -- Feind oder Verbündeter?}

Dies ist der tiefste konzeptuelle Unterschied zwischen beiden Dokumenten:

\begin{table}[h]
    \centering
    \caption{Die gegensätzliche Rolle von Rauschen in beiden Frameworks}
    \label{tab:noise_role}
    \resizebox{\textwidth}{!}{\begin{tabular}{@{}lll@{}}
        \toprule
        Aspekt & Dok.\ 034: QM-Optimierung & Dok.\ 161: T0-Ising Machine \\
        \midrule
        Rauschen ist... & \textbf{Feind} -- stört Gate-Treue & \textbf{Werkzeug} -- ermöglicht Annealing \\
        T0-Ursache & Fraktale Dämpfung $\Kfrak$ & Thermische Fluktuationen im OPO \\
        Strategie & Aktiv unterdrücken & Kontrolliert einsetzen \\
        Mechanismus & Zeitfeld-Pumpe moduliert $T(x)$ & Temperaturparameter $T_{\text{ann}}$ \\
        Optimum & Rauschen $\rightarrow 0$ & Rauschen $\rightarrow$ Minimum (dann $\rightarrow 0$) \\
        T0-Analogon & $\Kfrak \rightarrow 1$ (perfekte Gatter) & $\Kfrak^{3/2}$ als Annealing-Kraft \\
        \bottomrule
    \end{tabular}}
\end{table}

\begin{important}[Auflösung des Widerspruchs]
    Der scheinbare Widerspruch -- Rauschen ist in Dok.\ 034 schlecht, in Dok.\ 161 gut --
    löst sich durch die Skalierungsebene auf: Auf Gate-Ebene ist jedes Rauschen schädlich.
    Auf Systemebene ist kontrolliertes Rauschen notwendig, um lokale Minima zu verlassen.
    T0 beschreibt denselben physikalischen Mechanismus ($\Kfrak$) auf beiden Ebenen.
\end{important}

\subsection{Parallele 2: Fraktale Dämpfung als universeller T0-Mechanismus}

\begin{table}[h]
    \centering
    \caption{$\Kfrak$ als universeller T0-Mechanismus auf verschiedenen Ebenen}
    \label{tab:kfrak_universal}
    \resizebox{\textwidth}{!}{\begin{tabular}{@{}llll@{}}
        \toprule
        Kontext & Dok.\ 034 & Dok.\ 161 & Physikalische Bedeutung \\
        \midrule
        $\Kfrak$-Definition & $1 - 100\xipar$ & $K_{\text{frak}}^{3/2}$ (g-2-Korrektur) & Fraktale Raumzeit-Struktur \\
        Wirkung auf... & X-Gate-Rotation (dämpft) & Annealing-Dynamik (reguliert) & Abweichung vom Ideal \\
        Mathematisch & $\alpha = \pi \cdot \Kfrak$ & $\mathcal{D}(N) = e^{-\xipar \ln N/\Dfrak}$ & Geometrische Dämpfung \\
        Messbar durch & Residualer Flip-Fehler & $\ln(N)$-Residuum nach CIM & $\xipar$-Signaturen \\
        Testbar an & IBM-Quantencomputer & CIM mit variabler Spingröße & Beide Plattformen \\
        \bottomrule
    \end{tabular}}
\end{table}

\subsection{Parallele 3: Frequenzoptimierung vs.\ Topologieoptimierung}

Beide Dokumente schlagen Optimierungen der \emph{Systemarchitektur} vor -- auf
verschiedenen Ebenen:

\begin{table}[h]
    \centering
    \caption{Architekturoptimierungen: Gate-Level vs.\ System-Level}
    \label{tab:architecture}
    \resizebox{\textwidth}{!}{\begin{tabular}{@{}lll@{}}
        \toprule
        Optimierungstyp & Dok.\ 034: QM-Optimierung & Dok.\ 161: T0-Ising Machine \\
        \midrule
        Frequenzoptimierung & Harmonische Resonanzfrequenzen $f_n$ & -- \\
        Topologieoptimierung & T0-Topologie-Compiler (fraktale Weglänge) & 4D-Torus-Topologie \\
        Dekohärenzschutz & Aktive Zeitfeld-Modulation & Toroidale Randbedingungen \\
        Kopplungsstruktur & Gate-Pulse vorkompensiert & $J_{ij}$ geometrisch emergiert \\
        Ziel & Max.\ Gate-Treue & Min.\ Ising-Hamiltonian \\
        \bottomrule
    \end{tabular}}
\end{table}

\subsection{Parallele 4: Determinismus und Messbarkeit}

\begin{table}[h]
    \centering
    \caption{Determinismus und testbare Vorhersagen beider Dokumente}
    \label{tab:determinism}
    \resizebox{\textwidth}{!}{\begin{tabular}{@{}lll@{}}
        \toprule
        Aspekt & Dok.\ 034: QM-Optimierung & Dok.\ 161: T0-Ising Machine \\
        \midrule
        Art des Determinismus & Explizit: geometrische Transformationen & Versteckt: OPO-Feldgleichungen \\
        Quelle der Stochastik & Fraktale Raumzeit-Fluktuationen & Quantenvakuum-Anfangsbedingungen \\
        Testbare Vorhersage 1 & Gedämpfter X-Flip: $\alpha < \pi$ & $\ln(N)$-skalierendes Residuum \\
        Testbare Vorhersage 2 & Sweet Spots bei 6,24/2,38 GHz & $J_{ij}$-Abweichungen vom freien Param. \\
        Plattform & IBM Brisbane/Sherbrooke & CIM mit 1000+ Spins \\
        Quantifizierung & $\Delta\alpha = \pi(1 - \Kfrak) \approx 10^{-2}$ & $\xipar \approx 1{,}3 \times 10^{-4}$ \\
        \bottomrule
    \end{tabular}}
\end{table}

\clearpage

\section{Neue Optimierungsmöglichkeiten durch die Verbindung}
\label{sec:new_opportunities}

Die systematische Gegenüberstellung beider Dokumente eröffnet Optimierungsmöglichkeiten,
die in keinem der beiden Dokumente allein sichtbar waren:

\subsection{Hybride Optimierung: Gate-Ebene + Systemebene}

\begin{keyresult}[Hybride T0-Optimierungsstrategie]
    Eine vollständige T0-optimierte Quantenarchitektur kombiniert beide Ebenen:
    \begin{enumerate}
        \item \textbf{Gate-Ebene} (Dok.\ 034): Harmonische Frequenzkalibration +
        Zeitfeld-Pumpe für maximale Gate-Treue
        \item \textbf{Systemebene} (Dok.\ 161): T0-Ising Machine für kombinatorische
        Optimierung mit geometrisch emergenten Kopplungen
        \item \textbf{Verbindungsebene}: Der T0-Topologie-Compiler (Dok.\ 034) minimiert
        fraktale Weglängen -- dies ist dasselbe geometrische Prinzip wie die
        Torusstruktur in Dok.\ 161
    \end{enumerate}
\end{keyresult}

\subsection{Rauschen als Brücke zwischen beiden Ebenen}

Die dualistische Rauschrolle eröffnet eine neue Optimierungsstrategie:

\begin{itemize}
    \item \textbf{Phase 1 (Annealing)}: Kontrolliertes Rauschen ($\Kfrak$-Regime)
    -- die T0-Ising Machine sucht globale Minima
    \item \textbf{Phase 2 (Präzision)}: Rauschunterdrückung (Zeitfeld-Pumpe)
    -- Quantengatter führen die Lösung mit maximaler Treue aus
    \item \textbf{Übergang}: Der fraktale Parameter $\Kfrak$ reguliert beide Phasen --
    als Annealing-Mechanismus (Dok.\ 161) und als Gate-Dämpfung (Dok.\ 034)
\end{itemize}

\begin{equation}
    \text{Optimierung} = \underbrace{\Kfrak^{3/2}\text{-Annealing}}_{\text{Dok.\ 161:\ globales Minimum}}
    + \underbrace{\Kfrak \rightarrow 1\text{-Präzision}}_{\text{Dok.\ 034:\ Gate-Treue}}
    \label{eq:hybrid_optimization}
\end{equation}

\subsection{Neue testbare Vorhersage: Korrelation zwischen Gate-Fehler und CIM-Residuum}

Wenn $\Kfrak$ denselben universellen T0-Parameter auf beiden Ebenen beschreibt, dann
gilt:

\begin{equation}
    \frac{\Delta\alpha_{\text{Gate}}}{\pi} = \frac{\mathcal{D}_{\text{CIM}}(N)}{\ln N}
    = \xipar + O(\xipar^2)
    \label{eq:universal_kfrak}
\end{equation}

Das ist eine \textbf{neue quantitative Vorhersage}: Der gemessene Gate-Flip-Fehler an
IBM-Quantencomputern sollte mit dem $\ln(N)$-Residuum in CIM-Messungen
\emph{auf denselben numerischen Wert} $\xipar \approx 1{,}3 \times 10^{-4}$ konvergieren.
Dies ist expermenell überprüfbar mit heutiger Technologie.

\clearpage

\section{Zusammenfassung: Komplementäre Ebenen desselben Prinzips}
\label{sec:summary}

\begin{table}[h]
    \centering
    \caption{Vollständige Synthese: Beide Dokumente als komplementäre T0-Ebenen}
    \label{tab:synthesis}
    \resizebox{\textwidth}{!}{\begin{tabular}{@{}llll@{}}
        \toprule
        Dimension & Dok.\ 034 & Dok.\ 161 & Verbindung \\
        \midrule
        Ebene & Gate / Qubit & System / Spins & Selber $\xipar$ \\
        Rauschen & Feind (unterdrücken) & Werkzeug (kontrollieren) & $\Kfrak$ auf beiden Ebenen \\
        Darstellung & Zylindrisch $(z,r,\theta)$ & Phase $0°/180°$ & $T \cdot m = 1$-Zustände \\
        Topologie & Fraktale Weglänge & 4D-Torus & Geometrische Minimierung \\
        Kopplung & Gate-Pulse (vorkompensiert) & $J_{ij}$ (emergiert) & Zeitfeld $\Tfield$ \\
        Vorhersage & $\Delta\alpha \approx 10^{-2}$ & $\ln(N)$-Residuum & $\xipar = 4/30000$ \\
        \bottomrule
    \end{tabular}}
\end{table}

Die zentrale Einsicht lautet:

\begin{revolutionary}[Die vereinigte T0-Optimierungsperspektive]
    Sowohl Gate-Rauschen (Dok.\ 034) als auch kombinatorische NP-Optimierung
    (Dok.\ 161) sind Manifestationen \emph{desselben geometrischen Prinzips}:
    physikalische Relaxation in das Minimum des T0-Energiefunktionals
    $E_{\text{T0}} = \int |\nabla \Tfield|^2 \, d^4x$.
    Der Unterschied ist nicht konzeptuell -- er ist die Skalierungsebene der
    Betrachtung. Die fraktale Korrekturgröße $\Kfrak$ erscheint auf beiden Ebenen
    als derselbe universelle Parameter mit numerisch übereinstimmenden Vorhersagen.
\end{revolutionary}

\begin{thebibliography}{9}

\bibitem{pascher:qmopt}
J.~Pascher,
\newblock T0-Quantenmechanik: Geometrischer Formalismus und Optimierungsstrategien
für Quantencomputer,
\newblock T0-Dokumentationsreihe, Dokument 034 (2025).

\bibitem{pascher:ising}
J.~Pascher,
\newblock T0-Theorie und Coherent Ising Machines: Strukturelle Parallelen,
\newblock T0-Dokumentationsreihe, Dokument 161 (2026).

\bibitem{pascher:fundamentals}
J.~Pascher,
\newblock T0-Theorie: Zeit-Masse-Dualität als fundamentales Prinzip,
\newblock T0-Dokumentationsreihe, Dokument 003 (2025).

\bibitem{pascher:g2}
J.~Pascher,
\newblock Fraktale Korrekturen und die g-2-Anomalie,
\newblock T0-Dokumentationsreihe, Dokument 018 (2025).

\bibitem{pascher:quantum}
J.~Pascher,
\newblock T0-Quantencomputing: Qubits und Logikgatter,
\newblock T0-Dokumentationsreihe, Dokument 147 (2025).

\bibitem{pascher:scramblons}
J.~Pascher,
\newblock T0-Projektion auf Scramblon-Physik,
\newblock T0-Dokumentationsreihe, Dokument 148 (2026).

\bibitem{inagaki_cim_2016}
T.~Inagaki et al.,
\newblock A coherent Ising machine for 2000-node optimization problems,
\newblock \emph{Science} \textbf{354}, 603 (2016).

\end{thebibliography}
